% !TEX program = lualatex
\documentclass[12pt,a4paper]{article}

% --- Sprache und Schriftarten ---
\usepackage{fontspec}
\usepackage{microtype}
\usepackage[ngerman]{babel}

\IfFontExistsTF{CaskaydiaCove Nerd Font Propo}
  {\setmainfont{CaskaydiaCove Nerd Font Propo}}
  {\setmainfont{Latin Modern Roman}}

% --- Mathematik und Symbole ---
\usepackage{amsmath}
\usepackage{amssymb}
\usepackage{bm}
\usepackage{nicefrac}

% --- Layout und Formatierung ---
\usepackage[left=3.5cm,right=3.5cm,top=3cm,bottom=3cm]{geometry}
\setlength{\parindent}{0.0em}
\allowdisplaybreaks

% --- Grafiken und Links ---
\usepackage{graphicx}
\usepackage[hidelinks]{hyperref}
\usepackage{bookmark}

% --- Kopf- und Fußzeile ---
\usepackage{fancyhdr}
\pagestyle{fancy}
\fancyhf{}
\fancyfoot[C]{\thepage}
\renewcommand{\headrulewidth}{0pt}
\renewcommand{\footrulewidth}{0pt}

\newcommand{\titleline}{Diskrete Strukturen - Hausaufgabe [XX]}
\newcommand{\subtitleline}{06.1 Konstantin Di 8 Uhr MA 544}
\newcommand{\authorsline}{Sami Ertug Akkoca: [Matrikelnummer], Fabian Aps:525528, Max Kosec: [Matrikelnummer], Franz Juri Münzner: [Matrikelnummer]}
\renewcommand{\headrulewidth}{0pt}
\renewcommand{\footrulewidth}{0pt}

% --- Eigene Befehle und Umgebungen ---
\newcommand*{\QED}{\hfill\ensuremath{\square}}

\newenvironment{aufgabe}[1]{\noindent\textbf{#1.\,Aufgabe}\\[.5ex]}{\ \\[1ex]}
\newenvironment{loesung}[1]{\noindent\textbf{Lösung zu #1.\,Aufgabe}\\[.5ex]}{\vfill}


\hypersetup{
  pdftitle={\titleline},
  pdfauthor={\authorsline},
  pdfsubject={\titleline - \subtitleline},
  pdfcreator={LuaLaTeX}
}
% =======================================================
\begin{document}

% --- Manuelle Titelseite ---
\begin{titlepage}
	\centering
  {\Large Technische Universität Berlin\par}
  \vspace{2cm}
  {\Huge\bfseries \titleline\par}
  \vspace{1.5cm}
  {\Large \subtitleline\par}
  \vfill
  {\large Autoren: \authorsline\par}
  \vspace{2cm}
  {\large \today\par}
\end{titlepage}

% --- Inhaltsverzeichnis ---
\tableofcontents
\newpage

% --- Aufgabenseite ---
\section*{Aufgaben}
\addcontentsline{toc}{section}{Aufgaben}

\begin{aufgabe}{2}
	Es sei $x\in\mathbb{R}$. Zeigen Sie, dass für alle $n\in\mathbb{N}$ gilt, dass $x^n\in\mathbb{R}$.
\end{aufgabe}

\vspace{1cm}

\begin{aufgabe}{3}
	Es sei $f:\mathbb{R} \to \mathbb{R}, x\mapsto x^2\cdot e^{-x}$. Bestimmen Sie die erste Ableitung von $f$.
\end{aufgabe}

\newpage

% --- Lösungsseite ---
\section*{Lösungen}
\addcontentsline{toc}{section}{Lösungen}

\begin{loesung}{2}
	\textbf{Induktionsanfang:}
	Nach Voraussetzung ist $x^1=x\in\mathbb{R}$. Damit gilt schon einmal der Induktionsanfang.\\[0.5ex]

    \textit{Induktionsvoraussetzung:}
	Sei nun $x^n\in\mathbb{R}$ für ein beliebiges aber festes $n\in\mathbb{N}$.\\[.5ex]

    \underline{Induktionsschritt $(n\to n+1)$}: Ist die Induktionsvoraussetzung für $n\in\mathbb{N}$ erfüllt, so gilt, dass
	\[
	x^{n+1}
	= x^n\cdot x
	\in \mathbb{R}
	\]
	sein muss, da sowohl $x^n\in\mathbb{R}$ als auch $x\in\mathbb{R}$ und $\mathbb{R}$ abgeschlossen ist unter Multiplikation.
	\QED
\end{loesung}

\newpage

\begin{loesung}{3}
	Beide Abbildungen $x\mapsto x^2$ und $x\mapsto e^{-x}$ sind differenzierbar. Daher folgt nach der Produktregel:
	\begin{align*}
		f'(x)
		&= (x^2)'\cdot e^{-x} + x^2\cdot(e^{-x})'\\[2.5ex]
		&= 2xe^{-x} + x^2(-e^{-x}) \\
		&= (2x-x^2)e^{-x}.
	\end{align*}
	\QED
\end{loesung}

% Doom Emacs Lokale Variablen
%%% Local Variables:
%%% mode: latex
%%% TeX-engine: luatex
%%% TeX-master: t
%%% End:
\end{document}
